\documentclass[12pt]{article}
\usepackage[utf8]{inputenc}
\usepackage[T1]{fontenc}
\usepackage{amsmath,amssymb}
\usepackage{geometry}
\geometry{margin=2.5cm}

\begin{document}

\noindent $n\in\mathbb{N}$.

\bigskip
\noindent $H \overset{\text{not}}{=} n\mathbb{Z} = \{ny \mid y\in\mathbb{Z}\}$ is a subgroup of the group $(\mathbb{Z},+)$.

\noindent $\Rightarrow \exists!\, n\in\mathbb{N}$ such that $H=n\mathbb{Z}$.

\begin{enumerate}
\item[1)] $(\forall)\, x,y\in n\mathbb{Z} \Rightarrow x+y\in n\mathbb{Z}$.
\item[2)] $(\forall)\, x\in n\mathbb{Z} \Rightarrow -x\in n\mathbb{Z}$.
\end{enumerate}

\[
\left.\begin{aligned}
x&=nq_1\\
y&=nq_2
\end{aligned}\right\} \Rightarrow x+y = n(q_1+q_2).
\]

\bigskip
\noindent\underline{Converse.}

\noindent $H$ is a subgroup of $(\mathbb{Z},+)$ $\Rightarrow$ $H=n\mathbb{Z}$.

\noindent $0\in H$

\noindent $H=\{0\}$: \quad $\{0\} = \{0y \mid y\in\mathbb{Z}\}$.

\noindent Suppose $H\neq\{0\}$.

\[
H^+ = \{x\in H \mid x>0\}.
\]

\noindent Let $n$ be the smallest element of $H^+$

\noindent We show that $H=n\mathbb{Z}$.

\bigskip
\noindent (by induction)\footnote{the exact word/symbol is unclear in the original -- it appears to indicate an inductive argument.} $x=nq$.

\noindent $q=1 \Rightarrow x=n$

\noindent $nk\in H \Rightarrow n(k+1)\in H \Rightarrow nk+n\in H$.

\noindent Hence $n\mathbb{Z}\subseteq H$

\bigskip
\noindent Suppose\footnote{``Pp.'' could also mean ``by the proposition [of division with remainder]''.} $x=nq+r$.

\noindent Suppose $r\neq 0 \Rightarrow r=x-nq\in H$.

\bigskip
\noindent\underline{\textbf{Left (right) congruences on a group.}}

\noindent\underline{\textbf{Lagrange's Theorem.}}

\bigskip
\noindent Let $G$ be a group, $H\subseteq G$

\[
R_H^L = \{(x,y)\in G\times G \mid x^{-1}y\in H\}.
\]
\[
R_H^R = \{(x,y)\in G\times G \mid xy^{-1}\in H\}.
\]

\end{document}
